\subsection{The septic proper quotient}

The same CM arithmetic supplies the additional proper quotient needed
for the local statement at $3$.
Put $C_7:v^2+v=u^7$ and $A_7=\operatorname{Jac}(C_7)$.
\begin{proposition}\label{prop:septic}
The group $A_7(\mathbb Q)$ has order seven, and
$\Sha(\mathbb Q,A_7)$ is finite. In particular the $3$-adic Selmer group
of $T_3A_7$ is zero.
\end{proposition}
\begin{proof}
Use $(r,s,t,\delta)=(1,1,5,1)$ in Shu's Jacobi-sum description
\citep[\S\S2.1--2.2, pp.~5--6; Proposition 3.12(2), p.~16;
Theorem 3.13, pp.~17--18]{ShuExpanded2021}.
For $a=(-2)^5$, the integer $(a^6-1)/7$ is congruent to $3$ modulo
$7$. The Hecke character over $K=\mathbb Q(\zeta_7)$ therefore has
conductor $(1-\zeta_7)^2$. Its conductor over $\mathbb Q$ is $7^7$.
The finite root number at seven is $-\left(15/7\right)=-1$, while
its archimedean root number is $\left(-1/7\right)=-1$; all other
local signs are one. Thus the global sign is $+1$, and the completed
function is
\[
 \Lambda_7(s)=7^{7s/2}(2\pi)^{-3s}\Gamma(s)^3L(A_7,s),
 \qquad \Lambda_7(s)=\Lambda_7(2-s).
\]
These are algebraic Hecke-function inputs, not conclusions drawn from
numerical evaluations.

Here are the precise data and inequalities for central nonvanishing.
Write $B=7^{7/2}/(2\pi)^3$ and $L(A_7,s)=\sum a_nn^{-s}$.
The Mellin-integral derivation used for the degree-eleven quotient,
with three gamma factors, gives
\[
 L(A_7,1)=2\sum_{n\ge1}\frac{a_n}{n}Q_3(n/B),\qquad
 Q_3(x)=\frac1{2\pi i}\int_{(2)}
                   \frac{\Gamma(1+z)^3}{z}x^{-z}\,dz.
\]
Equivalently $Q_3(x)$ is the probability that the product of three
independent exponential random variables of mean one exceeds $x$.
For every positive integer $k$,
$Q_3(x)\le\Gamma(k+1)^3x^{-k}$ by Markov's inequality.
The Jacobi-sum Euler product gives
$|a_n|\le \sqrt n\,b_K(n)$, where
$\zeta_K(s)=\sum b_K(n)n^{-s}$.
Thus for $\sigma>1$ and $k+1/2>\sigma$ the omitted tail after $N$ is
at most
\[
 2\Gamma(k+1)^3 B^k N^{\sigma-k-1/2}\zeta_K(\sigma).
\]

All coefficients through $N=4096$ are computed exactly. At a prime
$p\ne7$ of even multiplicative order $f$ modulo seven, the polynomial
is $(1+p^{f/2}T^f)^{6/f}$, by
\citet[Lemma 1.1, p.~203]{GrossRohrlich1978}.
At residue degree three only $p=2,11$ require nonzero coefficients
below this cutoff. Direct counts over $\mathbb F_{p^3}$ give
\[
 P_2(T)=1-2T^3+8T^6,\qquad
 P_{11}(T)=1-44T^3+1331T^6.
\]
The first count is three points over $\mathbb F_8$; the second is
$1200$ points over $\mathbb F_{1331}$. At split primes the program
counts the seven character-value classes of the Jacobi sum, and
Newton identities give the required initial coefficients of its norm
polynomial. The factor at seven is one. Multiplicativity gives the
entire coefficient list; it has $124$ nonzero terms. For example,
$P_{29}(T)=1+T-27T^2+\cdots$.

For completeness, on the line $z=2+it$ the finite central sum is
\[
 \frac2\pi\operatorname{Re}\int_0^{20}
 \frac{\Gamma(3+it)^3}{2+it} B^{2+it}
 \sum_{n\le4096}\frac{a_n}{n^{3+it}}\,dt
\]
up to its vertical tail. Iterating the gamma recurrence seventeen
times and using the integral bound on the numerator gamma function
bounds that tail by
\[
 \frac2\pi B^2\left(\sum_{n\le4096}\frac{|a_n|}{n^3}\right)
 \frac{\Gamma(20)^3}{51\cdot20^{51}}.
\]
Arb arithmetic at $112$ bits encloses the displayed finite integral
by $0.97638934177248\pm5.84\cdot10^{-15}$ and the vertical tail by
$1.342\cdot10^{-16}$. With $k=4$ and $\sigma=11/10$, the Dirichlet
tail is less than $9.049\cdot10^{-6}$; the required zeta value is
$\zeta_K(11/10)=3.49720060785743449\ldots$, evaluated as the product
of the Riemann zeta factor and the five nonprincipal Dirichlet
$L$-functions of conductor seven. Consequently
\[
                 0.9763<L(A_7,1)<0.9765.
\]

Over the totally real field $F=\mathbb Q(\zeta_7+\zeta_7^{-1})$,
Shu's construction realizes the three two-dimensional coefficient
components of $A_7$ as parallel-weight-two Hilbert forms with trivial
central character: the restriction of the Hecke character to the real
idele group is the quadratic character of $K/F$, so automorphic
induction has central character equal to its square
\citep[\S3.4, p.~15]{ShuExpanded2021}. Each conjugate Jacobi Hecke function has the same
central value, since the conjugation reindexes the prime ideals of
$K$. The degree $[F:\mathbb Q]=3$ is odd, so
\citet[Theorem B and B-prime, p.~593; \S2.10.5, p.~641]{Nekovar2012}
applies without the alternative nonprincipal-local-component condition.
As in the preceding arithmetic arguments, it gives finite Mordell--Weil
group and finite whole Tate--Shafarevich group over $F$; restriction
then gives both finiteness assertions over $\mathbb Q$.
Finally \citet[Theorem 1.1, p.~202]{GrossRohrlich1978} identifies the
rational torsion as $\mathbb Z/7\mathbb Z$. Its exceptional septic
case does not occur: $1^3=1$ while $5^3=-1$ modulo seven.
Hence the rational group has order seven. Its $3$-adic completion
and the $3$-adic Tate module of its finite Tate--Shafarevich group both
vanish, proving the last assertion.
\end{proof}
